\begin{ouverture}[calcul]{6}{Découvrir les fractions}{Partie 2 · Fractions et nombres décimaux}
{Parts égales · Numérateur · Au-delà de l’unité}
\objectif{Lire, écrire et représenter une fraction.}
\objectif{Écrire une fraction supérieure à 1 avec un entier et une fraction.}
\objectif{Encadrer une fraction entre deux entiers consécutifs.}
\prerequis{Partages équitables et division.}
\end{ouverture}
\pagestyle{otMain}\essentiel\fontsize{11}{15}\selectfont
\phantomsection\label{cours-6-1}
\essSection{1}{Des parts égales}
\begin{definition}Quand une unité est partagée en 8 parts égales, chaque part est un huitième, noté $\frac18$. La fraction $\frac38$ désigne 3 de ces parts. Le nombre du bas, le dénominateur, indique le partage ; celui du haut, le numérateur, compte les parts prises.\end{definition}
\label{fig-cours-6-1-cours-parts}\figureGeo{cours-parts}

\begin{exemple}Dans une tablette de 20 carreaux, chaque carreau est $\frac1{20}$ de la tablette. Dire « les trois quarts », c’est prendre 3 parts d’un partage en 4.\end{exemple}
\begin{attention}Une fraction exige des parts égales. Si une figure est coupée en trois morceaux de tailles différentes, le morceau colorié n’est pas forcément $\frac13$ de la figure. Les parts peuvent avoir des formes différentes, à condition d’avoir la même taille.\end{attention}
\clearpage
\phantomsection\label{cours-6-2}
\essSection{2}{Des fractions plus grandes que 1}
\begin{propriete}Une fraction peut dépasser l’unité : $\frac74$, c’est 7 quarts, soit une unité entière (4 quarts) et encore 3 quarts. On écrit $\frac74=1+\frac34$. Si le numérateur dépasse le dénominateur, la fraction dépasse 1.\end{propriete}
\begin{exemple}$\frac93=3$ : neuf tiers remplissent exactement trois unités. $\frac{13}5=2+\frac35$ : avec 13 cinquièmes, on remplit 2 unités et il reste 3 cinquièmes.\end{exemple}
\begin{savoirfaire}[Encadrer entre deux entiers]Pour encadrer $\frac{11}4$ entre deux entiers : $\frac84=2$ et $\frac{12}4=3$, donc $2<\frac{11}4<3$. On cherche les paquets complets de 4 quarts.\end{savoirfaire}
\begin{savoirfaire}[Partager avec une fraction]Partager 7 galettes entre 4 enfants : chaque galette coupée en 4 donne 28 quarts, soit 7 quarts chacun. Chaque enfant reçoit $\frac74$ de galette, c’est-à-dire $1+\frac34$. La fraction écrit le résultat exact du partage.\end{savoirfaire}

\clearpage\rubrique{Découverte}
\begin{blocLivre}{ex-6-1}\begin{exercice}[1]{Lire une part}
Quelle fraction du disque la partie colorée représente-t-elle ?
\label{fig-ex-6-1}\figureGeo{frac-1}
\end{exercice}\end{blocLivre}
\begin{blocLivre}{ex-6-2}\begin{exercice}[1]{Colorier une fraction}
Reproduire la bande et colorier $\frac25$ de sa surface.
\label{fig-ex-6-2}\figureGeo{frac-2}
\end{exercice}\end{blocLivre}
\begin{blocLivre}{ex-6-3}\begin{exercice}[1]{Le bon vocabulaire}
Pour la fraction $\frac7{10}$, donner le numérateur et le dénominateur, puis écrire en chiffres « trois quarts » et « huit sixièmes ».
\end{exercice}\end{blocLivre}
\begin{blocLivre}{ex-6-4}\begin{exercice}[1]{Des parts vraiment égales ?}
La figure est coupée en trois morceaux et un morceau est coloré. Peut-on dire que la partie colorée représente $\frac13$ de la figure ? Expliquer.
\label{fig-ex-6-4}\figureGeo{frac-4}
\end{exercice}\end{blocLivre}
\begin{blocLivre}{ex-6-5}\begin{exercice}[1]{Au-delà de l’unité}
Chaque disque est partagé en quarts. Écrire la fraction colorée, puis l’écrire avec un entier et une fraction plus petite que 1.
\label{fig-ex-6-5}\figureGeo{frac-5}
\end{exercice}\end{blocLivre}

\rubrique{Entraînement}
\begin{blocLivre}{ex-6-6}\begin{exercice}[2]{Entier caché}
Écrire $\frac74$ sous la forme d’un entier et d’une fraction inférieure à 1. Que vaut $\frac93$ ?
\end{exercice}\end{blocLivre}
\begin{blocLivre}{ex-6-7}\begin{exercice}[2]{Recomposer}
Écrire $2+\frac35$ sous la forme d’une seule fraction.
\end{exercice}\end{blocLivre}
\begin{blocLivre}{ex-6-8}\begin{exercice}[2]{Encadrer une fraction}
Encadrer $\frac{11}4$ puis $\frac{17}5$ entre deux entiers consécutifs.
\end{exercice}\end{blocLivre}
\begin{blocLivre}{ex-6-9}\begin{exercice}[2]{Sept galettes pour quatre}
Quatre enfants se partagent équitablement 7 galettes identiques. Quelle part de galette chacun reçoit-il ? L’écrire de deux façons.
\end{exercice}\end{blocLivre}
\begin{blocLivre}{ex-6-10}\begin{exercice}[2]{Égales à un entier}
Compléter : $\frac{\square}6=1$ ; $\frac{12}{\square}=1$ ; combien de quarts font 2 unités ?
\end{exercice}\end{blocLivre}

\rubrique{Problèmes}
\begin{blocLivre}{ex-6-11}\begin{exercice}[3]{Le gâteau de la fête}
Un gâteau est découpé en 12 parts égales. Les invités mangent 5 parts le midi et 4 parts le soir.\begin{questions}
\item Quelle fraction du gâteau a été mangée ?
\item Quelle fraction du gâteau reste-t-il ?
\end{questions}
\label{fig-ex-6-11}\figureGeo{frac-11}
\end{exercice}\end{blocLivre}
\begin{blocLivre}{ex-6-12}\begin{exercice}[3]{La frise de la classe}
Une frise compte 20 carreaux identiques. Les élèves peignent en bleu le quart de la frise. Combien de carreaux sont bleus ? Quelle fraction de la frise reste à peindre ?
\label{fig-ex-6-12}\figureGeo{frac-12}
\end{exercice}\end{blocLivre}
\begin{blocLivre}{ex-6-13}\begin{exercice}[3]{Des tartes en sixièmes}
Deux tartes identiques sont coupées en sixièmes. Les convives mangent 7 parts. Écrire la fraction de tarte mangée, la comparer à 1 et l’écrire avec un entier.
\end{exercice}\end{blocLivre}
\begin{blocLivre}{ex-6-14}\begin{exercice}[3]{Dépasser 1 sans calculer}
Sans poser de calcul, dire si chaque fraction est inférieure, égale ou supérieure à 1 : $\frac98$ ; $\frac56$ ; $\frac{12}{12}$ ; $\frac{20}{19}$.
\end{exercice}\end{blocLivre}
\begin{blocLivre}{ex-6-15}\begin{exercice}[3]{La moitié de qui ?}
Hugo dit : « Une demi-pizza, c’est toujours plus qu’un quart de pizza. » Léna répond : « Pas si les pizzas sont différentes ! » Qui a raison ? Expliquer.
\end{exercice}\end{blocLivre}

\clearpage
